\documentclass[preprint]{JASA}

\begin{document}

\title[JASA/Sample JASA Article]{Adaptation rate and persistence across multiple sets of spectral cues for sound localisation}

\author{Paul Friedrich}
\email{paul.friedrich@uni-leipzig.de}\affiliation{Department of Biology, Leipzig University, Leipzig, Germany}

\author{Marc Schönwiesner}
\affiliation{Department of Biology, Leipzig University, Leipzig, Germany}

\preprint{Paul Friedrich, JASA}		

\date{\today} 

\begin{abstract}
\input{../Abstract/abstract.tex}
\end{abstract}

\maketitle

\section{\label{sec:1} Introduction}
% Not written yet. The earmold paper's introduction (inline in the old master)
% is archived at ../_archive_earmold_20260826/Manuscript_original.tex.


% ---------------------------------------------------------------------------
% METHODS — section master file.
% \input by main/Manuscript.tex. Each \subsection lives in its own file in this
% folder; this file holds only the section header, the opening paragraph, and
% the running order. Add new subsections as files and \input them below.
%
% ORDER follows the earmold paper (Friedrich & Schoenwiesner 2025): behavioral
% methods first, acoustic methods after Statistical analysis.
%
% WARNING (2026-08-26): main/Manuscript.tex \inputs this file AND still carries
% a full inline copy of the old earmold Methods further down. Until that inline
% block is removed, the section renders TWICE and \label{sec:2} is defined
% twice. See the note at the bottom of this file.
% ---------------------------------------------------------------------------

\section{\label{sec:2} Methods}

Participants had their individual HRTF recorded and trained to localize sounds
with one ear only in a virtual auditory environment (VAE) for 4 consecutive
days. To capture the trajectories of participants' adaptation to the HRTFs
presented in the VAE, participants completed daily localization tests. During
training and testing, HRTFs were presented to one ear only.

% --- behavioral methods ----------------------------------------------------

\subsection{Participants}
% TODO — n, age range, ethics approval, compensation. Recruitment ran on
% 18-40 years / normal hearing / 4 consecutive days. Planned n = 8, raised to
% 16 if the statistics require it.

\subsection{Experimental procedure}
% TODO — the 4-day timeline. Day 1: native binaural reference, donor build,
% baseline A and baseline D. Adaptation days: PRE test, training, POST test.
% Final day: counterbalanced 2x2 (Ear x Side), Williams Latin square.

% ---------------------------------------------------------------------------
% DRAFT SEGMENT — experimental setup
% Standalone; not \input anywhere yet. Stitch into Methods later.
% Code: experiment/localization/Localization_AR.py   (rendering + response)
%       hrtf/binsim/hrtf2binsim.py                   (settings, reverb, hp filter)
%       hrtf/binsim/hrir2mat.py                      (level normalization)
% Room / array / response apparatus are cited to the earmold paper rather than
% re-described — they are unchanged, and every reported recording came through
% record_dome on the existing speaker table (params.txt, all subjects).
% ---------------------------------------------------------------------------

\subsection{Experimental setup}

Recordings, localization tests and training were carried out in the
hemi-anechoic room and 45-loudspeaker array described previously
% \citep{friedrich_adaptation_2025}
, using the same head-mounted laser and electromagnetic motion sensor
(MetaMotionRL, Mbientlab, San Francisco, USA) for visual feedback of head
orientation and for the response. Acoustic measurements used the vertical
midline column of seven loudspeakers, spanning $\pm37.5^\circ$ elevation at
$12.5^\circ$ spacing. Signals were generated and acquired with TDT System 3
hardware (Tucker Davis Technologies, Alachua, USA) at a sampling rate of
48.8\,kHz.

Virtual sound sources were rendered over headphones (DT 990 Pro,
beyerdynamic, Heilbronn, Germany) by real-time convolution in pyBinSim
% \citep{neidhardt_pybinsim_2017}
, which selected the impulse response pair for the listener's current head
orientation and crossfaded between successive filters. Convolution ran in
blocks of 256 samples on a 512-tap direct-sound filter, followed in series by a
late reverberation tail at a direct-to-reverberant ratio of 20\,dB and the
participant's own headphone equalization filter. The gain of the rendered chain is
not fixed by the measurement, so each database was scaled by a single
broadband factor to a common value; presentation
level was thereby made identical across participants, conditions and days,
while interaural and between-direction level differences, being ratios within
a database, were unaffected. Overall level was matched by ear to the
free-field array and the system volume held fixed thereafter.

% --- open items ------------------------------------------------------------
% - SELF-CITATION verified: Friedrich, P., and Schoenwiesner, M. (2025).
%   "Adaptation rate and persistence across multiple sets of spectral cues for
%   sound localization," J. Acoust. Soc. Am. 157(3), 1543-1553,
%   doi:10.1121/10.0036056. Key friedrich_adaptation_2025; still needs adding
%   to the .bib, which does not yet exist for this manuscript.
% - pyBinSim citation not verified. Neidhardt et al., pyBinSim, is the usual
%   one; check what the pfriedrich-hub/pybinsim_tuil fork should be cited as,
%   and whether the fork's deviations need a sentence.
% - The late reverberation tail: where does it come from? Generic or measured?
%   One clause is owed, a reviewer will ask whether it is individual.
% - End-to-end latency of the rendering chain is NOT stated. One block is
%   5.24 ms at 48.8 kHz, but total motion-to-sound latency includes the sensor
%   rate and buffering and has not been measured. Either measure it and give a
%   number, or say nothing -- do not quote the block period as the latency.
% - Headphone model: DT 990 Pro confirmed for all reported participants
%   (MYSPHERE appears in the code as an alternative, unused here).
% - LEVEL JUSTIFICATION REWORDED 2026-08-26. It used to read "impulse responses
%   derived from directional transfer functions carry no absolute scale", which
%   rested on the DTF framing. A properly referenced HRTF DOES have a defined
%   magnitude (relative to free field), so that justification no longer holds.
%   The 6.5-15.2 dB chain-gain spread across databases is therefore NOT
%   explained by "no absolute scale" -- it must come from the processing
%   (windowing, onset alignment, low-frequency extrapolation) or from a real
%   between-subject difference, and 8.7 dB is too large to be the latter.
%   WORTH CHECKING: if the HRTFs are properly referenced, where does the spread
%   come from? The current wording is deliberately neutral and does not claim a
%   cause. Related: project_binsim_level_normalization.
% - Room dimensions deliberately not repeated; if a reviewer wants the paper
%   standalone, restore "6.4 x 4.4 x 2.3 m" to the first sentence.


\subsection{Localization task}
% TODO — sector grid (7 deg az, 14 deg el), el +-35, 3 targets per sector,
% min separation 20 deg, midline excluded from one-sided blocks; 75 trials per
% one-sided block, 150 for the binaural full-field reference. Also the
% free-field control: same task as Friedrich & Schoenwiesner (2025) but
% restricted to the seven midline loudspeakers, 3 targets each (21 trials),
% because the midline is the only direction at which the virtual rendering
% reproduces the measurement rather than the spherical-head model.

% ---------------------------------------------------------------------------
% DRAFT SEGMENT — localization test stimuli
% Standalone; not \input anywhere yet. Stitch into Methods later.
% Backing analysis + parameter sweeps: docs/stimulus_spectral_variation.md
% Code: experiment/localization/localization_helpers/stimulus.py
% ---------------------------------------------------------------------------

\subsection{Localization stimuli}

Localization stimuli were 225\,ms trains of five 25\,ms pink-noise bursts
separated by 25\,ms silent gaps (10\,ms raised-cosine onset and offset ramps,
5\,ms at the gap edges), cut from a single continuous noise realization. In the
fixed-spectrum condition the train was used unmodified, so its expected
spectrum was the same on every trial and only the realization noise varied
(SD of the 1/6-octave-smoothed log spectrum: 0.4\,dB across trains, 0.8\,dB
across the bursts within a train).

In the variable-spectrum condition each train was recolored, before gapping, by
a random spectral envelope drawn anew on every trial and shared by all five
bursts. Natural sound sources differ from one another mainly in broad spectral
coloration, which the auditory system must discount in order to read the finer,
elevation-dependent pattern imposed by the pinna; the envelope was built to
impose exactly that kind of variation. It was a smooth random function of log
frequency (0.5--16\,kHz), band-limited to ripple densities below 0.5 ripples
per octave with an rms of 3\,dB (median peak-to-trough depth 10\,dB), applied as
a zero-phase gain. The cutoff follows Macpherson and Middlebrooks
% \citet{macpherson_vertical-plane_2003}
, who found that localization was disrupted by spectral ripple at densities of
0.5--2 ripples per octave but not by coarser variation, and concluded that
broad-scale spectral structure is discounted despite contributing most to the
shape of directional transfer functions. Each trial was
therefore strongly recolored while the 0.5--2 ripples/octave band, in which the
pinna notch shifts with elevation, was left largely untouched. The manipulation
thus tests whether what was learned during training is tied to the spectrum of
the training stimulus or generalizes across source spectra. It does not
separate a perceptual remapping from an explicitly learned association: a
listener who had learned to read the preserved band would succeed under either
account. Timing, bandwidth and level were unchanged, and the envelope of every
trial was stored.

Across 400 realizations the variable stimulus varied with an SD of 2.4\,dB,
against 0.4\,dB for the fixed-spectrum stimulus. Within the 0.5--2
ripples/octave band the elevation-dependent variation of the participants'
head-related transfer functions (SD across elevation at fixed azimuth) exceeded
the trial-to-trial variation of the source spectrum by a factor of 7.5; below
0.5 ripples/octave the two were comparable (factor 1.0).

% NB: every number in this file scales linearly with the envelope rms, so
% changing the depth is a proportional edit throughout. Values at the earlier
% 8 dB setting, for reference: 27 dB median peak-to-trough, 6.4 dB stimulus SD,
% 2.8 in the cue band, 0.4 below it. Settle this against the self-test
% (protocols/dev/stimulus_check.py cells 7-9) before submission.

\subsubsection*{Choice of envelope depth}

Envelope depth is bounded from both sides, and the two bounds are what fixed
it. It must be large enough that the coarse-scale coloration of the source is
no longer a reliable guide to elevation, and small enough that the elevation
cue itself survives. The upper bound is the more concrete of the two.
Macpherson and Middlebrooks
% \citet{macpherson_vertical-plane_2003}
report degradation of vertical-plane localization at 1 ripple per octave only
for peak-to-trough depths at or above 20\,dB, and at coarser densities even
40\,dB left performance largely intact; we therefore treated 20\,dB
peak-to-trough as a ceiling and stayed well below it. Two further
considerations argue for the same direction. The band limit is not sharp---a
function band-limited on a finite log-frequency axis has skirts---so a deeper
envelope leaks into the 0.5--2 ripples/octave region even though no energy is
assigned there, and the depth of the elevation cue itself varies between ears
and listeners, so a setting chosen against one head-related transfer function
must leave headroom for shallower ones.

The lower bound is set by the purpose of the manipulation. In the fixed-spectrum
condition the sound at the eardrum carries a near one-to-one map from absolute
spectrum to elevation, so a listener can in principle solve the task by
matching the overall spectrum to a stored template, without separating source
from filter. The envelope must make that route unreliable, which requires its
trial-to-trial variation to be at least comparable to the elevation-dependent
variation of the head-related transfer function at the same spectral scales. At
3\,dB rms the two are of the same magnitude below 0.5 ripples per octave
(factor 1.0), so overall coloration no longer specifies elevation, while the
cue band retains a 7.5:1 margin. Deeper envelopes would make the source
dominate rather than merely match at coarse scales, but only at the cost of
the margin in the cue band, and of exceeding the depth at which disruption has
been reported.

Variation was confined to the coarse band rather than spread across all
spectral scales for the same reason: variation inside the cue band would
contaminate the very structure being read, and contributes nothing to defeating
a template-matching strategy that the coarse band does not already achieve. It
was retained as an adjustable parameter and set to zero throughout, so that
deliberately contaminating the cue band remains available as a manipulation
check.

Finally, the parameters were verified behaviorally rather than only
acoustically. [PLACEHOLDER — self-test, cells 7-9: own unmodified HRTF,
binaural, fixed-spectrum vs variable-spectrum blocks; report elevation gain for
each and the difference. With intact ears and an unmodified transfer function,
a drop in elevation gain can only reflect the envelope encroaching on the cue,
so an absent or small drop is the direct evidence that the envelope is
transparent to localization.]

% --- supplement ------------------------------------------------------------
% The envelope is a discrete cosine series in dB on a 128-point log-frequency
% grid, coefficient k corresponding to k/10 ripples per octave; the constant
% term and all coefficients at or above 0.5 ripples/octave are zero, a zero
% coefficient contributing 0 dB, i.e. unity gain. Burst-to-burst correlation of
% the detrended log spectra within a trial was r = 0.95, versus r = -0.01 for
% the unmodified stimulus.

% --- what this manipulation can and cannot decide ---------------------------
% Rules out: a mapping tied to the overall spectrum of the training stimulus
%   ("this timbre = that elevation"), since gross coloration now varies trial to
%   trial by as much as elevation itself varies it. Note the strength of this
%   claim depends on the depth: at 3 dB the source only MATCHES the cue at coarse
%   scales, it does not swamp it, so the timbre route is made unreliable rather
%   than uninformative. Phrase it that way in the Discussion; "uninformative"
%   would need the deeper envelope, which M&M's depth threshold rules out.
% Does NOT rule out: a lookup restricted to the preserved 0.5-2 rip/oct band.
%   The pinna notch is a shape-within-band feature, so a remapping of that
%   structure onto space and a learned association between that structure and a
%   response both survive the manipulation. Arguably the two are not
%   distinguishable acoustically at all -- adaptation at the spectral-to-spatial
%   mapping stage IS the learning of new pattern-location associations.
% What would separate perceptual from explicit/procedural learning:
%   - response latency and dual-task cost (explicit strategies are slow and
%     interference-prone)
%   - generalization to untrained locations WITHIN the trained ear: a graded map
%     should interpolate, a lookup should not (cf. condition B)
%   - graded vs catastrophic degradation under band-limiting of the cue region
%   - aftereffects / coexistence on return to the native HRTF
%   - subjective report: does it sound like it comes from there, or is the
%     response reasoned to
% On whether the listener CAN discount the source spectrum here — two reasons
% they can, so a failure is informative rather than a ceiling of the paradigm:
%  (1) Presentation is not strictly monaural. The other ear receives the same
%      source convolved with a cepstrally smoothed (n_keep=4) DTF, so the
%      interaural log-difference cancels the source spectrum exactly while
%      retaining almost all of the trained ear's fine structure. The reference
%      channel is in fact smoother than in natural binaural listening, where
%      both ears carry notches. Its level falls with lateral angle (up to ~9 dB
%      down at 35 deg), so the reference degrades toward the periphery.
%  (2) Source-spectrum discounting is anyway achievable monaurally: models based
%      on the first/second spectral derivative recover direction for sources
%      with locally flat or locally constant-slope spectra (Zakarauskas &
%      Cynader 1993), and current sagittal-plane models use monaural spectral
%      gradients (Baumgartner et al. 2014, already cited in the Intro).
%
% Framing for the Discussion: Macpherson & Middlebrooks show the NORMAL system
% already discounts broad-scale spectral variation. The open question is whether
% a NEWLY acquired spectral-to-spatial mapping inherits that invariance, or is
% initially stimulus-specific and only becomes source-invariant with further
% exposure. That, not "sensory vs cognitive", is what the ripple blocks test.

% --- open items ------------------------------------------------------------
% - Citations to add once a .bib exists:
%     Macpherson & Middlebrooks (2003) JASA 114(1):430-445,
%       doi:10.1121/1.1582174  -- ripple-spectrum probe, the 0.5-2 rip/oct window
%     Zakarauskas & Cynader (1993) JASA 94(3):1323-1331 -- spectral-derivative
%       model, monaural source-spectrum invariance
%     Middlebrooks (1992), Baumgartner et al. (2014) -- already cited in Intro
% - DONE: depth is now reported peak-to-trough as well as rms, and the depth
%   bound is argued explicitly ("Choice of envelope depth").
% - Cue depth varies by ear and subject; the band ratios are scaled from FS's
%   modified left ear at az -20. Decide whether to report a group range instead
%   -- with a shallow envelope the cue-band margin is comfortable everywhere, so
%   a group minimum would be the more convincing number to give.
% - The self-test paragraph is a placeholder until cells 7-9 have been run.
%   If the fixed vs variable difference in elevation gain is not small, the
%   depth argument above has to be rewritten around whatever value survives.
% - FS's two "USO" blocks were run with the old generator (one fixed base
%   texture per session, 0.8 dB variation). They are NOT a variable-spectrum
%   test and should be reported as such, or dropped.


% ---------------------------------------------------------------------------
% DRAFT SEGMENT — the cue manipulation and how the donor is chosen
% Standalone; not \input anywhere yet. Stitch into Methods later.
% Code: hrtf/modify/donor_detail.py       (the composite)
%       hrtf/processing/midline.py        (applied to the measured arc)
%       hrtf/analysis/donor_selection.py  (vanopstal_dtfs / isim / composite_isim)
%
% Code and text now agree: shortlist()/select_donor() gate on ridge_slope <= 0.5
% and rank by |r_match - TARGET_R_MATCH|, both from pair_metrics() on the detail
% domain; DONOR_POOL holds the behaviourally qualified donors and
% MIN_DONOR_ELEVATION_GAIN is the criterion. isim()/vanopstal_dtfs() remain in
% the module as diagnostics and must NOT be reported as similarity.
% ---------------------------------------------------------------------------

\subsection{Spectral cue manipulation}

Participants were trained and tested with a modified version of their own
head-related transfer functions in which the coarse spectral envelope was their
own and the fine spectral detail was another listener's. For each direction and
ear, the log-magnitude spectrum was separated by cepstral truncation into an
envelope, retaining the first four coefficients, and a residual containing
everything finer. The modified spectrum was the participant's own envelope plus
a donor's residual. Own phase was retained and each direction was rescaled to
its original broadband energy, so interaural time differences and broadband
interaural level differences were unchanged and only spectral shape finer than
the envelope scale was altered.

Retaining four coefficients follows Kulkarni and Colburn
% \citet{kulkarni_role_1998}
, who smoothed transfer functions by cepstral truncation and found listeners
could no longer discriminate smoothed from original once roughly sixteen
coefficients were retained; four is well inside the range where the change is
audible and leaves the envelope too coarse to carry an elevation-dependent
pattern. Splitting the spectrum this way rather than substituting the donor's
transfer function whole was verified to decorrelate the spectral-to-spatial map
as thoroughly as a whole substitution, whereas the complementary
composite---donor envelope with own detail---left the native map intact. Detail
therefore carries the elevation cue and the envelope carries timbre, allowing
the cue to be replaced while level, head shadow and individual coloration are
retained.

The manipulation was applied to the measured median-plane transfer functions
and the full directional set was regenerated from them through the
spherical-head model used to build the unmodified set (Sec.~\ref{sec:hrir}), so
that native and modified sets share their interaural time differences exactly
and the envelope is fitted to the pinna response rather than to the pinna
response multiplied by the head shadow.

\subsection{Donor qualification}

Donor transfer functions were taken from a pool of previously recorded
participants. A recording qualified as a donor only if that participant had
themselves localized accurately with their own unmodified transfer function in
the same virtual environment, quantified as the elevation gain of their first
completed localization block (criterion: gain $\geq 0.8$; targets within
$\pm30^\circ$ elevation). The manipulation transplants one listener's spectral
detail to another, and this criterion establishes that the pattern being
transplanted is one a human auditory system demonstrably reads: it is evidence
that the cue works, obtained from the only instrument that can supply it.

The criterion is behavioral because the acoustic alternatives are not
predictive. Among the recordings for which both were available, the depth and
elevation dependence of a donor's own spectral detail---the natural acoustic
index of how much cue it carries---was unrelated to how well that participant
localized with it (Spearman $\rho = +0.18$, $n = 24$); the recording with the
most pronounced fine structure of the entire set belonged to a participant whose
elevation gain was $0.05$. Qualification is applied as an inclusion criterion
only. Accurate localization shows that a pattern is usable; inaccurate
localization is ambiguous between a weak cue, a rendering failure and an
inattentive listener, so an unqualified recording is treated as unproven rather
than as unsuitable.

\subsection{Donor selection}

Among qualified donors, the donor assigned to a participant was the one whose
composite came closest to an intended perturbation size while remaining
unabsorbable by that participant's existing spectral-to-spatial map. Both
quantities were computed between the participant's own and the modified spectral
detail on the median-plane arc, in the 5657--11314\,Hz band, averaged over ears.

Perturbation size was quantified as $r_\mathrm{match}$, the mean over elevations
of the correlation between the participant's own detail at that elevation and
the modified detail at the same elevation. This is the quantity that the
mechanism proposed by Van Wanrooij and Van Opstal
% \citet{vanwanrooij_relearning_2005}
---a correlation between the sensory input and a stored spectral representation
of the listener's own ears---predicts should govern performance. It was targeted
mid-range rather than minimized, because the perturbation is bounded from both
sides: in their monaural experiment, listeners whose cues were not sufficiently
perturbed never reached stable performance, while Trapeau et al.
% \citet{trapeau_fast_2016}
found that larger perturbations produced a greater initial loss and slower, less
complete recovery. The target was set to the median $r_\mathrm{match}$ between
qualified donors, so that the perturbation corresponds to a typical difference
between two people whose own cue is known to work.

Absorbability was quantified as the slope of the regression of the
best-matching own elevation on true elevation. A slope near unity means the
existing map can still read the modified set as a coherent elevation, displaced
rather than replaced, and can accommodate the manipulation as a constant
response bias; this is the configuration that Van Wanrooij and Van Opstal's
non-adapting listeners occupied. Candidates with a slope above $0.5$ were
excluded, and among the remainder the donor closest to the target
$r_\mathrm{match}$ was selected.

Both measures are required. Across all donor--listener pairings available here
they were nearly independent (Spearman $\rho = +0.14$), and of the pairings
whose $r_\mathrm{match}$ indicated a strong perturbation ($\leq 0.3$), $31\%$
nonetheless had a slope above the exclusion threshold---a large change in the
spectral cue that the existing map could still absorb.

Both were computed on the spectral detail rather than on whole transfer
functions. Correlating transfer functions requires the component common to all
directions to be removed by a reference derived from directions other than those
being compared, which for the spectral variability index
% \citep{trapeau_fast_2016}
is a diffuse-field average that the present recordings do not sample. A
reference formed from the median-plane arc itself cannot substitute, since it is
the mean of the same spectra that are then correlated: the residuals sum to zero
by construction, the mean pairwise correlation is fixed at $-1/(N-1)$, and the
resulting index is a constant (measured: $1.02 \pm 0.02$ across 31 recordings,
against $1.056$ predicted for $N = 19$ elevations). The cepstral residual is
subject to no such constraint, because the envelope removed from each direction
is a smoothed copy of that same direction and directions remain independent.

% --- open items ------------------------------------------------------------
% CITATIONS — verified 2026-08-18 unless marked:
%   Van Wanrooij & Van Opstal (2005) J Neurosci 25(22):5413-5424,
%     doi:10.1523/JNEUROSCI.0850-05.2005. "Relearning sound localization with a
%     new ear" (SINGULAR -- the MONAURAL perturbation study). Do NOT confuse
%     with Hofman, Van Riswick & Van Opstal (1998) Nat Neurosci, "with new ears"
%     (plural), the BINAURAL mold study. I_sim definition = their Materials and
%     Methods, "DTF similarity index"; DTF definition = their Eq. 1; the 0.5 and
%     0.3 anchors = their Fig. 7A and 7C; r^2 = 0.94 = their Fig. 6 (n=6).
%     All quoted from the full text, verified.
%   Kulkarni & Colburn (1998) Nature -- cepstral smoothing, discrimination
%     threshold near 16 coefficients. [volume/pages NOT VERIFIED]
%   Trapeau et al. (2016) JASA -- VSI. Cited only if the paragraph declining to
%     use it names it. [volume/pages NOT VERIFIED]
%
% - QUALIFIED POOL = 7 (CO, AGV, SW, VD, AH, FS, FD), from 27 recordings with an
%   own-HRTF block. Every one of the 12 listeners resolves to an admissible
%   donor: 6 in band, 6 widened, no fallbacks. Target r_match 0.58 = median over
%   the 21 qualified-donor pairs (min 0.32, Q1 0.51, Q3 0.66, max 0.77).
%   RECOMPUTE THE TARGET WHENEVER THE POOL CHANGES -- it is defined from the pool.
% - 0.8 falls in the gap between FD at 0.82 and AS at 0.63, so it is not making a
%   fine distinction. State it as a cut in a gap, not as a graded criterion.
% - GS IS A PROBLEM CASE: the only admissible qualified donor for GS is AGV at
%   r_match 0.90, i.e. almost no perturbation. Either GS needs a donor that is
%   not yet qualified, or GS cannot be run under this rule. Check before
%   scheduling GS.
% - Five candidates still have no own-HRTF block (GLK, IM, MSc, NW, UG) -- they
%   have other runs, just none with their unmodified HRTF. One ~60-trial block
%   each would qualify them and widen the pool.
% - CORRECTED: an earlier version of this file called the EG distribution
%   bimodal with a suspicious cluster at zero. That came from
%   vsi_rmse.find_participants() missing 6 subjects (it searched only the flat
%   results/pilot/<id>.pkl and not the nested results/pilot/<id>/<id>.pkl, and
%   preferred the pilot pair where an id existed in both trees, which shadowed
%   the active AS and PF). Both are fixed; with the fuller set the spread is
%   closer to continuous. Do not repeat the bimodality claim.
% - WHY NO I_SIM. Van Opstal's index was tested and rejected for this role: it is
%   the SD down each column of the correlation matrix, and SD is invariant to
%   permutation of the rows, so it cannot distinguish a listener's own ear from
%   their own ear relabelled by 17 deg (measured: identical value 0.542 for
%   shifts of 0, 4, 8 and 17 deg, while the correlation at matching elevation
%   fell 1.00 -> 0.08). It also scores a cue-REMOVED ear (own envelope only,
%   0.395) as less perturbed than a real donor composite (0.272), so it does not
%   separate replacement from removal. Keep this in the file; it is the answer if
%   a reviewer asks why the obvious measure was not used.
% - Detail strength has no repeatability estimate (every arc is measured once),
%   so do not report it to more precision than the ranking requires, and do not
%   claim a difference between adjacent donors.
% - NO REPEATABILITY ESTIMATE exists for any of these measures: every arc is
%   measured once. Re-recording one subject would fix this cheaply, and two
%   recordings HAVE been remade (AS 2026-08-18, GS 2026-07-28) -- but as
%   replacements, not repeats, so they do not serve.
% - Sec.~\ref{sec:hrir} is a placeholder label -- point it at whichever section
%   describes the HRIR recording and the spherical-head azimuth expansion.
% - Regenerate if anything changes: the own-ear and between-listener I_sim
%   values and their n, and the target once set.


\subsection{Statistical analysis}
% TODO — polar error as the primary metric, elevation gain / RMSE / residual
% SD secondary. Transfer measured as within-ear pre-post CHANGE SCORES, not
% the post-only A-vs-D contrast (the mirrored filter carries a foreign-base
% disadvantage that cancels in the change score). Report simple effects of Ear
% within laterality rather than interpreting the Ear x Location interaction.

% --- acoustic methods ------------------------------------------------------

% ---------------------------------------------------------------------------
% DRAFT SEGMENT — HRTF measurement and the virtual HRTF set
% Standalone; not \input anywhere yet. Stitch into Methods later.
% Two \subsection{} blocks. Placed LATE in Methods, where the earmold paper put
% "Binaural recordings" and "Directional transfer functions".
% Code: hrtf/record/recordings.py::record_dome        (acquisition, head tilts)
%       hrtf/record/record_hrir.py                    (orchestration)
%       hrtf/record/processing.py                     (deconvolve, equalize,
%                                                      lowfreq, expand_azimuths)
%       hrtf/record/fit_head_radius.py                (acoustic radius fit)
%       hrtf/record/calibration/calibrate_headphones.py
% Parameters below are read off the code and the recorded params.txt, not the
% module-level interactive defaults (which differ -- head_radius in particular).
% \label{sec:hrir} resolves the placeholder reference in
% methods_donor_selection.tex.
% ---------------------------------------------------------------------------

\subsection{HRTF measurement}\label{sec:hrir}

Head-related impulse responses were measured at the blocked entrance of each
ear canal with MEMS microphones (AMM-3738-T, PUI Devices, Dayton, USA) mounted
in anthropometric earplugs of the PIRATE design
% \citep{denk_pirate_2019}
, fitted to each participant from a size series
% \citep{bau_adapting_2024}
, and supplied with phantom power (PP2B, Millenium). Logarithmic sweeps of 200\,ms duration (120\,Hz to 22\,kHz, 85\,dB SPL) were
presented ten times from each loudspeaker of the midline column and averaged
after alignment. Seven loudspeakers at $12.5^\circ$ spacing sample elevation
too coarsely to resolve the spectral cue, so the arc was recorded three times,
the participant being cued to tilt the head by one third of the loudspeaker
spacing between passes. Interleaving the physical positions in this way yields
19 elevations from $-37.5^\circ$ to $37.5^\circ$ in steps of $4.17^\circ$. A reference recording was then made with the microphones mounted
on a stand at the position the head had occupied and with no listener present,
using identical settings; it captures the response of loudspeakers,
microphones and room in situ. Emitted sweeps were pre-filtered per loudspeaker
with inverse filters obtained from a flat probe microphone, which equalizes
each loudspeaker in level and magnitude response rather than merely matching
them to one another.

Impulse responses were recovered by regularized inversion of the sweep and
each was divided by the reference response of the same loudspeaker. The
quotient is the head-related transfer function (HRTF) in the usual sense: the
pressure at the blocked ear canal entrance relative to the free-field pressure
at the position of the head with the listener absent. These are HRTFs and not
directional transfer functions: no direction-independent component has been
removed. Our earlier work
% \citep{friedrich_adaptation_2025}
reported directional transfer functions, obtained by dividing each transfer
function by the average over all measured transfer functions. That step is
omitted here because only the midline is measured, so an average over the
measured set would be an average over the very directions that are
subsequently compared, and would remove part of the elevation dependence it is
meant to preserve. Each quotient was onset-aligned, windowed with a 2.5\,ms
Hann window and truncated to 512 samples.

\subsection{Virtual HRTF set}

Because only the midline arc was measured, the remaining directions were
synthesized with a rigid-sphere model of the head
% \citep{duda_range_1998}
. The radius of the sphere was fitted for each participant to interaural time
differences measured from a horizontal row of loudspeakers, taken from the
slope of the interaural phase between 200 and 1500\,Hz. The measured arc was then replicated
across azimuth in 25 steps of $4.17^\circ$ from $-50^\circ$ to $50^\circ$,
giving 475 directions, each carrying the spherical-head magnitude shading and
interaural phase difference relative to the frontal direction at the same
elevation. Below 800\,Hz, where the measurement is limited by the sweep
bandwidth and the room, magnitude was replaced by spherical-head values.

One consequence of this construction should be stated plainly. Spectral detail
in the delivered set is measured on the midline and is identical at every
azimuth; azimuth is carried by broadband interaural level and time differences
alone, and the rendering is acoustically faithful to the measurement only at
$0^\circ$ azimuth. Three considerations make this acceptable here. The
elevation-dependent spectral pattern, which is what the manipulation alters
and what the analysis measures, is measured rather than modeled. Every
condition in the design is rendered through the identical expansion, so no
comparison rests on the synthesized component. And because a magnitude-only
edit cannot move a modeled interaural time difference, the native and modified
sets share their interaural cues exactly, which would not be guaranteed had
both been measured. The design accordingly contrasts hemifields rather than
azimuthal acuity, and no claim is made about localization in azimuth.

Two residual corrections were applied. Division by the reference leaves a
small interaural difference at frontal directions, where none should exist;
this was removed by equating broadband energy between the ears over 200\,Hz to
16\,kHz, one scalar per direction, which on the manikin left a frontal
interaural level difference of $0.000 \pm 0.003$\,dB. Headphone equalization
was measured for each participant from three re-seatings of the headphones,
averaged, and inverted as a regularized minimum-phase filter of 256 taps.

% --- open items ------------------------------------------------------------
% - CITATIONS:
%     friedrich_adaptation_2025 -- VERIFIED: Friedrich, P., and Schoenwiesner,
%       M. (2025). J. Acoust. Soc. Am. 157(3), 1543-1553,
%       doi:10.1121/10.0036056.
%     denk_pirate_2019 -- VERIFIED (Zenodo record 2574395). NOTE THE AUTHOR
%       ORDER: Denk, F., Brinkmann, F., Stirnemann, S., and Kollmeier, B.
%       (2019). "The PIRATE: an anthropometric earPlug with exchangeable
%       microphones for Individual Reliable Acquisition of Transfer functions
%       at the Ear canal entrance," Fortschritte der Akustik -- DAGA, Rostock,
%       doi:10.5281/zenodo.2574395. Denk is first author, not Brinkmann.
%     bau_adapting_2024 -- VERIFIED: Bau, D., Moscher, O., and Poerschmann, C.
%       (2024). "Adapting the PIRATE Design for Cost-Effective In-Ear
%       Microphones," Fortschritte der Akustik -- DAGA, Hannover, pp.
%       1604-1606. This is the paper the present hardware follows; it tests
%       AMM-3738-T (PUI) alongside CMM-2718AT (CUI) and SPH1642-HT5H (Knowles)
%       and finds all three near-equivalent, so it supports the choice rather
%       than prescribing it. It is also the source for the XS-XL size series.
%     Duda & Martens (1998) JASA 104(5) -- "Range dependence of the response of
%       a spherical head model". CHECK this is the right one for the model as
%       implemented; if processing.py follows a different formulation, cite
%       that instead. [NOT VERIFIED]
% - MICROPHONE now stated: AMM-3738-T (PUI Devices) + PP2B phantom supply
%   (Millenium is Thomann's house brand -- CHECK whether to name Thomann as the
%   manufacturer instead, and confirm the town for the PUI address line).
%   CONFIRM ALSO that the plugs really are the XS-XL size series rather than
%   individually printed; "fitted to each participant from a size series" is
%   written on the strength of the setup following Bau et al. (2024).
%   The old PUM-3046L-R sentence is gone -- do not let it drift back in.
% - PER-SUBJECT RADII ARE NOT REPORTED. Give the group mean and range once the
%   cohort is closed. Also state what happens on a failed fit (the code falls
%   back to 0.0875 m and stores nothing) -- if any reported participant hit the
%   fallback, that must be said.
% - The head-radius paragraph is now bare by decision (Paul, 2026-08-26): no
%   manikin validation, no Woodworth comparison. If a reviewer asks why the
%   fitted radii (~0.072 m) fall below the Woodworth 0.0875 m, the answer is
%   that the low-frequency phase delay of a real head exceeds its onset delay
%   (Kuhn 1977) and 0.0875 m is the ONSET-delay standard, inappropriate for a
%   phase-based ITD. The fit uses the same estimator the model then imposes.
%   Keep this note; it is the rebuttal, held out of the text on purpose.
% - "limited by the sweep bandwidth and the room" is my reading of why 800 Hz;
%   confirm the actual reason before submission.
% - The 19 elevations come from 7 + 6 + 6 across the three tilt passes (the
%   outermost speaker drops out of the two tilted passes). Not stated in the
%   text -- add if a reviewer queries the count.
% - record_mesm (the newer 7-speaker path) was NOT used for any reported
%   subject. If later participants are recorded on it, this subsection needs a
%   variant, and the two cohorts cannot be pooled without checking.
% - Dome loudspeaker equalization was OFF for subject and reference alike
%   (equalize_dome: False in every params.txt), which is what makes the
%   reference division a clean in-situ calibration. Not stated; probably does
%   not need to be, but it is the answer if asked why no dome EQ is mentioned.


\subsection{Training}
% TODO — cite the training game from Friedrich & Schoenwiesner (2025) (pulse
% rate and inter-pulse delay as feedback, 90 s rounds, scoring, leaderboard)
% and describe only the delta: rendered over headphones through the VAE, one
% ear, and whatever changed in the region sampling.

% ---------------------------------------------------------------------------
% REMOVED 2026-08-26: the second paragraph of the old skeleton, which was still
% about mold 1 / mold 2 and five-day adaptation periods, i.e. the earmold
% paper. Kept here so nothing is lost:
%
%   Participants wore two distinct sets of binaural earmolds consecutively
%   (denoted as mold 1 and mold 2), each over the course of a five-day
%   adaptation period. The acoustical and behavioral impact of the molds were
%   measured, and the trajectories of participants' adaptation were recorded to
%   capture the occurrence of meta-adaptation. To test whether learning a new
%   set of spectral cues interferes with a previously learned mapping, the
%   persistence of adaptation to the first earmolds was measured after
%   participants adapted to the second molds. This measurement was repeated for
%   the second molds after five days without earmolds for comparison.
%
% The previous version of this file is saved as methods.tex.bak_20260826.
% ---------------------------------------------------------------------------


\section{\label{sec:3} Results}

\begin{figure}[t]
\includegraphics[width=17cm]{../figures/Figure4}
\caption{\label{fig:FIG4}{Time course of vertical localization performance. White circles indicate localization with free ears, black circles and grey squares indicate localization with the first and second earmolds respectively. Data points represent the mean across participants (mold 1: n = 15, mold 2: n = 12). Error bars show the standard error of the mean. Individual learning trajectories are shown in light grey. \textbf{A} Time course of elevation gain, \textbf{B} vertical localization error and \textbf{C} vertical response variability (SD) throughout the two consecutive adaptation periods. Participants elevation gain was reduced while localization error and variability were increased after insertion of earmolds. Elevation gain and vertical error recovered within each of the two five-day adaptation periods, while response variability further increased. Participants' gain recovered faster with the first molds than with the second set. After each mold removal, elevation gain returned to baseline levels whereas vertical error increased compared to the first test with free ears (white dots on day 5 and 10). Although elevation gain with the first molds was slightly decreased after participants adaptated to the second set, the persistence of adaptation did not differ between molds.}}
\end{figure}

\subsection{Earmolds reduced vertical localization performance}

\section{\label{sec:4} Discussion}
Earmolds repeatedly disrupted vertical localization, but participants relearned sound localization with both mold sets within successive five-day adaptation periods. Individual adaptation performance did not improve with the second earmolds and learning a second set did not interfere with the previously learned representation. Localization accuracy with free ears slightly decreased throughout the experiment, whereas elevation perception remained unaffected. Spectral information of free ears correlated with behavior, and was similarly reduced by both mold sets. Acoustic dissimilarities between successive pinna shapes (free ears vs. the first molds and the first vs. the second molds) were similar, though each set induced changes at distinct frequencies.

\subsection{Localization with free ears} 


\section{\label{sec:1} Introduction}
\subsection{Research topic}
Accurate localization of sounds is required for navigation, communication, predation and escape. It is therefore not surprising that neural and structural specializations have evolved to perform this complex task. For instance, the human cortex continuously integrates sensory input across modalities to estimate the relative direction and distance of objects in the environment. Whereas the topography of visual and somatosensory space is represented by a point‐by‐point mapping of primary receptors, the auditory scene must first be computed from incoming sound waves. In this scenario, the direction-dependent filter profiles of the pinnae and upper body  result in spectral cues that aid in the inference of sound source elevation and front-back discrimination (\citet{blauert_sound_1969},  \citet{middlebrooks_narrow-band_1992}, \citet{baumgartner_modeling_2014}). As these head related transfer functions change throughout life (\citet{keen_growth_1988}, \citet{otte_age-related_2013}), the auditory system must retain its ability to recalibrate the mapping of spectral cues to locations in space. 

\subsection{Previous work}
\citet{hofman_relearning_1998} demonstrated this recalibration experimentally by inserting molds in the concha of adult listeners. Localization accuracy was initially impaired by the binaurally modified spectral cues and recovered within a few weeks, as participants adapted to the earmolds. Interestingly, listeners' localization accuracy with their native ears remained unchanged once the molds were removed, indicating that accommodation to new spectral cues does not interfere with the representation of the original cues. The authors suggested that adapting to a new pinna shape resembles the learning of a new language, in that distinct sets of pinna filters are stored in parallel and without much interference. In this view, the auditory system selects the correct spectral-to-spatial map based on the current pinna filter. However, \citet{trapeau_fast_2016} showed that after adaptation, listeners were able to accurately localize sounds with their native ears at the first trial, without tactile information about the absence of the molds.  A lack of even short-lived aftereffects indicates that no switch between discrete maps is required and could be better explained by a many-to-one mapping mechanism, in which several spectral profiles become associated with one spatial location \citep{trapeau_fast_2016}. Early investigations suggested similar models, in which the spectra of incoming sounds are compared with stored spectral templates associated with different directions (e.g., \citet{blauert_sound_1969}; \citet{hebrank_spectral_1974}). Later studies provided further evidence supporting a spectral correlation model (\citet{middlebrooks_narrow-band_1992}, \citet{langendijk_contribution_2002}, \citet{baumgartner_modeling_2014}). While spectral cue learning resembles a small example of auditory plasticity, a deeper understanding of this model system could uncover more general principles of neural learning. 

\subsection{Current approach}
So far, studies investigating the adaptation process have focused on the effects of learning one additional set of spectral cues (see \citet{carlile_plastic_2014} for a comprehensive review). However, adaptation to multiple pinna shapes could reveal limitations in the capacity of the auditory system to simultaneously encode several discrete spectral-to-spatial mappings without interference and provide constraints for theoretical models of spectral cue encoding. In the present study, adult listeners learned to localize sounds with two different sets of earmolds within two consecutive five-day adaptation periods. To establish both binaural representations in quick succession, participants wore the molds continuously throughout the adaptation periods and underwent daily sessions of sensory-motor training, which has been shown to accelerate the adaptation to modified pinnae (\citet{carlile_relearning_2014}, \citet{parseihian_rapid_2012}, \citet{trapeau_fast_2016}). 

\subsection{Hypotheses}
This paradigm was selected to test two hypotheses: (1) Learning a new set of spectral cues is accelerated by plasticity induced during recent cue relearning, i.e. the occurrence of meta-adaptation (\citet{robinson_meta-adaptation_2016}, \citet{liu_higher-level_2020}). If this holds true, listeners should adapt faster to the second set of earmolds compared to the first set. (2) \citet{hofman_relearning_1998} demonstrated that learning new spectral cues does not interfere with the neural representation of listeners' original cues. Building on this, we hypothesize that this principle also applies to a recently learned representation, thus learning a second set of earmolds will not interfere with listeners' ability to localize sounds using the first set.

\section{\label{sec:2} Methods}

Participants wore two distinct sets of binaural earmolds consecutively (denoted as mold 1 and mold 2), each over the course of a five-day adaptation period. The acoustical and behavioral impact of the molds were measured, and the trajectories of participants’ adaptation were recorded to capture the occurrence of meta-adaptation. To test whether learning a new set of spectral cues interferes with a previously learned mapping, the persistence of adaptation to the first earmolds was measured after participants adapted to the second molds. This measurement was repeated for the second molds after five days without earmolds for comparison.

\subsection{Participants}
15 participants (11 females and 4 males age between 20 and 30 years) took part in the experiment. They were informed about the relevant experimental procedures before providing their consent. Participants had no history of hearing disorder and had normal or corrected vision. The experimental procedures were approved by the local ethics committee. Monetary compensation was provided to the participants based on the time they spend at the experiment.

\begin{figure}[t]
\includegraphics[width=12.6cm]{../figures/Figure1}
\caption{\label{fig:FIG1}{Time course of the experiment.}}
\end{figure}

\subsection{Experimental procedure}
During their first visit, participants were familiarized with the environment, equipment and procedures of the free-field localization task. To maximize procedural learning before the experiment so that any learning during the experiment could be attributed solely to sensory learning of the new spectral cues, participants completed at least one localization task and were free to continue practicing until they felt comfortable with the task. No feedback was given. After the initial familiarization, participants performed one localization task to measure their baseline localization accuracy. Once this task was completed, the directional components of participants' head related transfer functions, the directional transfer functions (DTFs), were acquired \citep{middlebrooks_individual_1999}. Participants’ ears were fitted with the first set of earmolds, after which they repeated the localization task. If elevation perception was not sufficiently disrupted, the molds were further adjusted until the subsequent localization test yielded an elevation gain below 0.4. To capture changes to spectral cues induced by the molds, DTFs were measured again with the molds in place. From that day on, participants continuously wore the silicone molds for five days without removing them. Skin adhesion of the fitted molds was checked during participants' daily visits to the lab. Throughout this adaptation period, participants underwent a daily routine of training sessions. Each training session was followed by a free-field localization task. On the final day of the first adaptation period (day 5 in Fig.~\ref{fig:FIG1}), participants completed a short training session and a localization test. To verify that adaptation generalized to new sounds, a subset of 6 participants performed an additional localization test presenting stimuli with varying spectral content (see section \ref{stimuli} for a detailed description of stimuli). The earmolds were then removed, and participants’ localization accuracy was immediately measured again to test for aftereffects of the adaptation. The second pair of earmolds was then fitted to the participants’ ears, their initial localization accuracy was measured and DTFs were acquired. Over the next 5 days, the procedure was repeated as described for the first adaptation period, including daily training sessions and localization tasks. After mold removal at the end of the second adaptation period (day 10 in ~\ref{fig:FIG1}), the first earmolds were briefly re-inserted and participants completed a localization run. To compare the persistence of adaptation to both earmolds, participants returned to the lab after 5 days for a final localization test with the second molds re-inserted. 

\subsection{Experimental setup}
The localization tests, binaural recordings and training sessions were conducted in a hemi-anechoic room (6.4 x 4.4 x 2.3 m). Participants were seated in a comfortable chair in front of a spheric array of 45 loudspeakers (Mod1, Orb Audio, New York, USA) with a radius of 1.4 m, covering the visual field. Loudspeakers were hidden by an acoustically transparent curtain to avoid visual cues of the sound source positions. A small light emitting diode that was visible through the curtain indicated the center of the frontal hemifield and served as fixation point (0° azimuth, 0° elevation). During the localization tasks and training sessions, participants wore a headband with a laser pointer and an electromagnetic motion sensor (MetamotionRL, Mbientlab, San Francisco, USA) attached. The laser light was reflected by the curtain and provided visual feedback for the participants to indicate sound source direction by pointing their head at the perceived location. Real time head orientation and position captured by the motion sensor were used to calculate the interaural-polar coordinates of participants’ responses.

\begin{figure}[t]
\includegraphics[width=14.5cm]{../figures/Figure2}
\caption{\label{fig:FIG2}{\textbf{A}, Example of a silicone mold in the right ear of a participant. \textbf{B-D} DTFs across 7 elevations on the vertical midline obtained from the right free and modified ear of a participant. Color maps show the spectral amplitude computed at 83 frequency bins between 4 and 16 kHz. The spectral profile was linearly interpolated between the measured DTFs for display. Acoustic measures computed for these DTFs: spectral information; free ears: 0.86, mold 1: 0.56, mold 2: 0.38, spectral dissimilarity; free ears and mold 1: 0.9, mold 1 and mold 2: 0.38, free ears and mold 2: 1.28.}}
\end{figure}

\subsection{Earmolds}
To alter participants’ perception of sound source elevation, their pinnae were modified by silicone molds. As a result, spectral cues derived from the shape of the external ears were changed sufficiently to diminish participants elevation discrimination. Earmolds were created by applying fast curing, skin safe and adherent silicone (SkinTite, Smooth-On, Macungie, USA) to the cymba conchae, cavum conchae and the antihelix while keeping the ear canal unobstructed (Fig.~\ref{fig:FIG2} A). To ensure similar levels of disruption in elevation perception across participants, the shape and volume of the earmolds were customized and iteratively adjusted. Three participants did not complete the experiment (mold 1: n = 15, mold 2: n = 12) and were removed from the analysis for these molds.

\subsection{Localization task} 
44 Loudspeakers were used for the localization task, covering 105° in azimuth (-52.5° to 52.5°) and 75° in elevation (-37.5° to 37.5°) of the hemisphere in front of the participant. Loudspeaker positions are described in an interaural-polar coordinate system. The loudspeakers were arranged on a spherical grid, formed by 7 loudspeakers on the horizontal and 7 loudspeakers on the vertical axis. Loudspeakers were distributed on the sphere with an angular distance of 17.5° in azimuth and 12.5° in elevation between neighboring speakers. The central and the four outermost loudspeakers (at $\pm$-52.5° azimuth and  $\pm$37.5° elevation) were excluded. At the beginning of each trial, participants were instructed to aim the head mounted laser at a centrally presented LED while pressing the button on a handheld box. When the button was pressed, the participant’s initial head orientation was recorded and the stimulus was presented at a pseudorandom direction. Participants were instructed to indicate the perceived direction by turning their head towards the sound source and to confirm their response by pressing the button again. The displacement in azimuth and elevation from the central position was defined as the localization response. No feedback was given. Each direction was presented three times during a localization run, with an angular distance of at least 35° between sound locations of consecutive trials to reduce neural adaptation.

\subsection{Stimuli} \label{stimuli}
The stimuli in the free-field localization task were 225 ms long trains of Gaussian broadband noise (20-22000 Hz). Stimuli consisted of five noise bursts of 25 ms duration with 1 ms sinusoid onset and offset ramps and 25 ms of silence between each burst. The stimuli were short to exclude head movement effects and were chosen following a study of \citet{dramas_auditoryguided_2008}, which showed an improvement of localization accuracy between repeated 40 ms noise bursts and a single 200 ms burst. Stimuli in the generalization test consisted of 225 ms long mixtures of unidentifyable environmental sounds. Each stimulus in this test was composed of 6 randomly arranged excerpts of short anechoically recorded impact sounds \citep{nakamura_acoustical_2000} drawn from a list of 42 recordings and had a unique spectrum. Stimuli were re-generated for each trial and controlled by the slab toolbox \citep{schonwiesner_soundlab_2021} in a custom Python script. Overall sound pressure level of the stimuli at the position of the participants’ ears was 42 dB. Stimuli were processed digitally and amplified via TDT System 3 hardware (Tucker Davis Technologies, Alachua, USA) at 50kHz sampling rate and 24bit amplitude resolution. To reduce differences in amplitude and frequency response across the loudspeakers, we applied an inverse finite impulse response filter for each speaker. To obtain these filters, frequency modulated sweeps presented from each loudspeaker were recorded by a probe microphone (Brüel \& Kjær, Nærum, Denmark) positioned at equal distance and orientation. Resulting transfer functions did not differ more than 3 dB between 1 and 20 kHz across equalized loudspeakers.

\subsection{Statistical analysis}
Sound localization was analyzed for accuracy, response variability and response gain. Each measure was quantified for elevation and azimuth components respectively. Accuracy was quantified by the localization error, defined as the root mean squared distance between the physical target and the perceived response locations \citep{hartmann_localization_1983}. Variability of participant’s responses was defined as the mean standard deviation (SD) across response angles for each sound location \citep{makous_two-dimensional_1990}. The participants ability to perceive sound source azimuth and elevation was quantified by the response gain, defined as the slope of the linear regression line between target and response angles \citep{hofman_relearning_1998} (see Fig. 1~\ref{fig:FIG3}). Throughout the analysis, paired comparisons were statistically assessed using Wilcoxon signed-rank tests. To compare three or more measurements the Friedman test was used for paired samples, and the Kruskal-Wallis test for independent samples. Relations between measures were determined using Spearman correlation coefficients. Participants who did not complete the experiment were removed for paired comparisons. The false positive rate was set to 5\%.

\begin{figure}[h]
\includegraphics[width=9.5cm]{../figures/Figure3}
\caption{\label{fig:FIG3}{Stimulus-response plot resulting from a localization test; \textbf{A} with free ears and \textbf{B} after mold insertion. Grey circles represent individual response elevations. Elevation gain was defined as the slope of the linear regression line (dashed) between target and response angles \citep{hofman_relearning_1998}. Black lines indicate target elevations, grey lines show the mean vertical error of individual responses for each target. Localization error (RMSE) was quantified by the mean absolute error across targets. Grey bars show the standard deviation of individual responses for a given target. Response variability (SD) was defined as the mean standard deviation across target angles.}}
\end{figure}

\subsection{Binaural recordings}
Binaural recordings were conducted to measure participant’s head related transfer functions with and without silicone molds. PUM-3046L-R miniature microphones (PUIaudio, Fairborn, USA) were inserted in the ear canal to measure the sound pressure level at the eardrum. To minimize non-directional contributions by standing-waves in the canal, the microphones were placed 2 mm into the entrance of the blocked ear canal. Participants were seated in front of the loudspeaker array and were asked to remain stationary during the measurement while aiming the head mounted laser at the central LED. Frequency modulated sweeps of 100 ms duration were presented 30 times from each of the loudspeakers on the vertical midline (from -37.5° to 37.5° elevation, at 0° azimuth). Recordings were averaged for every location to increase signal to noise ratio. Recordings were digitized via TDT system 3 hardware at a sampling rate of 97 kHz. Transfer functions were extracted from the time-averaged recordings by taking the ratio of the Fourier transform of the acquired signal to the Fourier transform of the input signal.

\subsection{Directional transfer functions}
To reduce non-directional portions of the head related transfer functions (such as ear canal resonance and microphone transfer function) each measured transfer function was divided by the average across transfer functions of all participants with and without molds (n = 588). DTFs were processed with a bank of logarithmically spaced cosine band-pass filters as proposed by \citet{middlebrooks_individual_1999}, reflecting the frequency dependent resolution of human hearing. Center frequencies were equally spaced at 0.0286 octaves between 4 and 16 kHz resulting in 2\% frequency difference between each of the 83 filters (Fig.~\ref{fig:FIG2} B-D). To quantify behaviorally relevant acoustic factors in the measured DTFs we used two measures previously introduced by \citet{trapeau_fast_2016}; vertical spectral information (VSI) and vertical spectral dissimilarity (VSI dissimilarity). Spectral information quantifies information in participants’ DTFs that could potentially guide elevation discrimination and was defined as one minus the mean of the coefficients obtained from the autocorrelation matrix of a set of DTFs (see Fig. 1 in \citet{trapeau_fast_2016}). A set of identical transfer functions for all elevations will result in a VSI of 0, whereas highly different transfer functions will result in a VSI close to 1. Spectral information previously correlated with vertical localization error and explained 25\% of the behavioral variance among individuals \citep{trapeau_fast_2016}. Spectral dissimilarity measures the amount of acoustical change in spectral cues caused by the pinna modifications. Spectral dissimilarity was defined as the root-mean-square distance between the correlation matrix of two sets of DTFs and the autocorrelation matrix of the DTFs measured with free ears. \citet{trapeau_fast_2016} showed that spectral dissimilarity correlated with loss in elevation performance following earmold insertion. Spectral information and spectral dissimilarity were computed from DTFs obtained from the left and right ear. For comparison with behavioral measures left and right ear values were averaged.

\subsection{Training}
The training task, inspired by \citet{trapeau_fast_2016}, was designed to accelerate the adaptation to new spectral cues and was introduced to participants as a game-like scenario. They were instructed to find the location of a pulsed pink noise played from one of the 44 speakers. Feedback was provided by varying pulse duration and the delay between the pulses depending on the participant’s head orientation. Duration and delay of consecutive pulses decreased logarithmically with decreasing angular distance between target sound direction and the listeners head orientation, from up to 500 ms at 65° angular distance. The pulse train gradually merged into a continuous sound when participants directed the head mounted laser to within 3° angular distance from the sound source. If participants remained oriented at the target area for at least 500 ms, the sound source was considered found. Participants scored points for every found location depending on how quickly it was found and a popular video game sound indicated the number of rewarded points. The target sound then switched to a new location at least 45° in joint angular distance (azimuth and elevation) away from the previous location. Participants tried to score as many points as possible within 90 seconds. The final score was displayed on a screen after each round and a leaderboard encouraged competition. Throughout the adaptation periods participants underwent a daily routine of three 10 minute training sessions, with 5 minute breaks between each session. 

\section{\label{sec:3} Results}

\begin{figure}[t]
\includegraphics[width=17cm]{../figures/Figure4}
\caption{\label{fig:FIG4}{Time course of vertical localization performance. White circles indicate localization with free ears, black circles and grey squares indicate localization with the first and second earmolds respectively. Data points represent the mean across participants (mold 1: n = 15, mold 2: n = 12). Error bars show the standard error of the mean. Individual learning trajectories are shown in light grey. \textbf{A} Time course of elevation gain, \textbf{B} vertical localization error and \textbf{C} vertical response variability (SD) throughout the two consecutive adaptation periods. Participants elevation gain was reduced while localization error and variability were increased after insertion of earmolds. Elevation gain and vertical error recovered within each of the two five-day adaptation periods, while response variability further increased. Participants' gain recovered faster with the first molds than with the second set. After each mold removal, elevation gain returned to baseline levels whereas vertical error increased compared to the first test with free ears (white dots on day 5 and 10). Although elevation gain with the first molds was slightly decreased after participants adaptated to the second set, the persistence of adaptation did not differ between molds.}}
\end{figure}

\subsection{Earmolds reduced vertical localization performance}
Insertion of earmolds acutely disrupted vertical sound localization (see Fig.~\ref{fig:FIG4}, day 0; free ears vs. mold 1: elevation gain: $W_{15} = 120, p < .001$, error: $W_{15} = 0, p < .001$, SD: $W_{15} = 17, p < .006$, day 5; free ears vs. mold 2: elevation gain: $W_{12} = 78, p < .001$, error: $W_{12} = 0, p < .001$, SD: $W_{12} = 0, p < .001$). Horizontal gain and accuracy were not affected by the earmolds, with azimuth gains of $\sim$ 1.0 and horizontal errors of $\sim$ 5.5° across measurements. However, horizontal response variability was slightly increased by the earmolds (free ears vs. mold 1: SD: $W_{15} = 28, p = .036$, effect size: 0.3°, free ears vs. mold 2: SD: $W_{12} = 17, p = .005$, effect size: 0.6°). Both sets of earmolds caused similar disruption of vertical localization (disruption on day 0 vs. day 5: elevation gain: $W_{12} = 25, p = 0.301$, error: $W_{12} = 29, p = 0.47$, SD: $W_{12} = 21, p = 0.176$). Compared to the first set, insertion of the second mold set appeared to impair localization less for sounds located below the midline (Fig.~\ref{fig:FIG8}).

\subsection{Participants adapted to both mold sets}
Participants wore the two sets of earmolds during two consecutive five-day adaptation periods. Adaptation was driven by their day-to-day multi-sensory experience and daily sessions of sensory-motor training at the lab. Vertical localization performance improved significantly for both molds except for response variability, which increased throughout the adaptation period (Fig.~\ref{fig:FIG4}, mold 1; day 0 vs. day 5: elevation gain: $W_{15} = 0, p < .001$, error: $W_{15} = 119, p < .001$, SD: $W_{15} = 15, p = .004$, mold 2; day 5 vs. day 10: elevation gain: $W_{12} = 0, p < .001$, error: $W_{12} = 63, p = .032$, SD: $W_{12} = 9, p = .008$). Horizontal localization was not affected by adaptation to the earmolds. 

\subsection{Adaptation generalized to new sounds}
To test whether participants memorized location-specific spectral features of the training stimuli for sound localization, the localization test was repeated with a subset of six participants on the last day of each adaptation period using unidentifiable environmental sounds of random spectral content. The effect of these stimuli on vertical localization error did not differ between adapted earmolds and free ears indicating that learning generalized to new sounds (Friedman test; differences in vertical error between pink noise and environmental sounds across conditions; free ears: $-0.38^\circ$ $\pm$ 1.82 (mean $\pm$ standard error), mold 1: $2.5^\circ$ $\pm$ 0.46, mold 2: $-0.15^\circ$ $\pm$ 0.45, $\chi^2_2 = 4, p = 0.135$).

\subsection{Adaptation rate did not increase with the second mold}
To investigate effects of meta-adaptation, i.e., faster relearning of sound localization with new pinna shapes after previously adapting to a different set of pinna modifications, the rate of adaptation was compared between the first and the second set of molds. We defined the rate of adaptation as the increase in vertical localization performance over a five day adaptation period. To account for acoustic differences between sets of earmolds, the rate of adaptation was divided by the initial decrease in localization performance caused by the respective mold set. Whereas elevation gain recovered faster with the first molds, the reduction of vertical error did not differ between mold sets (performance gain from day 0 to day 5 divided by initial performance reduction; mold 1 vs. mold 2: elevation gain: $W_{12} = 8, p = .006$). Adaptation rate varied continuously between participants and did not fall into discernible groups such as low and high performers. Individual differences in adaptation rate were not consistent across mold sets.

\subsection{Adaptation persistence did not differ between successive earmolds}
To determine whether learning a new set of spectral cues interfered with a previously learned set, localization with the first earmolds was tested again after participants adapted to the second set of molds. A final localization test with the second molds (day 15 in Fig.~\ref{fig:FIG4}) served as a control, as participants were not exposed to new spectral cues in the five days after the second molds were removed. Participants vertical error with the first mold set was slightly increased by the concurrent adaptation to the second set, whereas other measures were not affected (error with mold 1; day 5 vs. day 10: 17.12° $\pm$ 0.65 vs. 19.3° $\pm$ 1.17, $W_{15} = 22, p = .015$). Localization with the second molds remained unchanged after 5 days with free ears. Adaptation persistence was defined as the change in vertical localization performance between the last test with molds in the adaptation period and the final persistence test. No differences in adaptation persistence were found between the first and second mold set. Fig.~\ref{fig:FIG5} shows the average localization response throughout the experiment. 

 \begin{figure}[t]
\includegraphics[width=13cm]{../figures/Figure5}
\caption{\label{fig:FIG5}{Change of response pattern throughout the experiment. Grey panels show sound target locations, averaged for targets belonging to similar directions by dividing the target space into twenty-five half-overlapping sectors \citep{hofman_relearning_1998}. Black dots represent sound localization responses for each target averaged across all participants. The first two rows show behavioral impact, followed by adaptation after 5 days and adaptation persistence with mold 1 and 2, respectively.}}
\end{figure}

\subsection{Free ears localization accuracy decreased during adaptation}
In previous studies, participants' localization with free ears remained unaffected by adaptation to new spectral cues (\citet{hofman_relearning_1998}, \citet{trapeau_fast_2016}, \citet{carlile_relearning_2014}, \citet{van_wanrooij_relearning_2005}). To confirm these findings, free ears localization was tested immediately after mold removal at the end of each adaptation period. Whereas elevation gain and response variability remained unchanged, participants’ vertical localization accuracy was slightly reduced after each adaptation period (Fig.~\ref{fig:FIG4} B, error with free ears; day 0 vs. day 5: $W_{15} = 9, p = .001$, day 0 vs. day 10: $W_{12} = 3, p < .001$). 

\subsection{Spectral information increases localization performance}
Participants' DTFs were recorded across the vertical midline with and without earmolds. To quantify the information available for vertical sound localization in the individual sets of DTFs, vertical spectral information, i.e., the dissimilarity between DTFs across elevations, was computed in 5 octave bands between 4 and 16 kHz (4–8 kHz, 4.8–9.5 kHz, 5.7–11.3 kHz, 6.7–13.5kHz, 8–16kHz). Spectral information of participants’ free ears varied among frequency bands ($\chi^2_4 = 13.59, p = .009$), and peaked in the 5.7–11.3 kHz band as previously reported by \citet{trapeau_fast_2016} (Fig.~\ref{fig:FIG6} A). Spectral information measured in this band is shown for all participants in Fig.~\ref{fig:FIG6} B. Spectral information of participants' left and right ears was similar (correlation between spectral information of individuals left and right ears in the 5.7–11.3 kHz band: $R = 0.45, p = .0897$). To account for variations in the frequencies of spectral features among individuals reported by \citet{middlebrooks_individual_1999}, spectral information was additionally computed in a wider frequency range (5.7–13.5 kHz), where it correlated with vertical localization error (Fig.~\ref{fig:FIG6} C) and response variability (free ears vertical spectral information in the 5.7–13.5 kHz band compared to vertical localization: error: $R = -0.58, p = .025$, SD: $R = -0.58, p = .025$). Spectral information was therefore used to quantify information for vertical sound localization in the following analyses.

\begin{figure}[t]
\includegraphics[width=7cm]{../figures/Figure6}
\caption{\label{fig:FIG6}{Free ears vertical spectral information (VSI) and localization performance. \textbf{A} Variation of spectral information across five octave bands between 4 and 16kHz. Spectral information was averaged across sets of DTFs obtained from participants’ left and right ears without earmolds. Spectral information varies significantly and is highest in the 5.7–11.3 kHz band.  \textbf{B} Left and right ear spectral information of each participant in the 5.7-11.3 kHz band. Spectral information of the left and right ears were similar. \textbf{C} Relation between spectral information in the 5.7 – 13.5 kHz band and vertical localization accuracy.}}
\end{figure}

\subsection{Earmolds reduced spectral information in the 5.7–11.3 kHz band}
To verify whether earmolds induced similar acoustic changes, their acoustic properties were carefully compared. Application of silicone molds to the pinnae altered spectral cues in the 4–16 kHz band. The consecutively applied sets of earmolds attenuated the spectral notch situated in the 5–11 kHz band and the neighboring spectral peak in the 11–14 kHz band (Fig.~\ref{fig:FIG7}). Similar effects have been observed in previous studies using pinna modifications to alter spectral cues for sound localization (\citet{trapeau_fast_2016}, \citet{van_wanrooij_relearning_2005}, \citet{hofman_relearning_1998}). Comparing spectral information of participants’ modified and free ears in the 5.7–11.3 kHz band showed that earmolds reduced the information available for vertical localization in this band (see Fig.~\ref{fig:FIG8} B, differences between spectral information of free and modified ears; free ears vs. mold 1: $W_{30} = 325, p = .002$, free ears vs. mold 2: $W_{24} = 251, p = .001$). Both mold sets led to similar reduction of spectral information compared to participants’ free ears (reduction in spectral information caused by mold 1 vs. mold 2: $W_{24} = 102, p = 0.197$). The relation between spectral information of participants’ left and right ears persisted after mold insertion but was not significant for the second molds (mold 1: $R = 0.56, p = .034$, mold 2: $R = 0.53, p = .064$). Fig.~\ref{fig:FIG8} A shows the spectral information of all participants with modified and free ears. \citet{van_wanrooij_relearning_2005} reported the formation of new acoustic cues at higher frequencies after the insertion of silicone molds. To test whether silicone molds increased the frequencies of spectral features, spectral information was additionally compared in the 11.3–13.5 kHz band. The first molds increased spectral information in this band compared to free ears ($W_{30} = 125, p = .039$).

\begin{figure}[t]
\includegraphics[width=11.5cm]{../figures/Figure7}
\caption{\label{fig:FIG7}{Acoustic effect of the earmolds. \textbf{A} Mean across participants’ DTFs with free ears. \textbf{B-C} Probability maps showing the number of participants for which the molds induced marked changes (see section ~\ref{change_n}) in spectral amplitude at each elevation and frequency bin. Black contour lines include elevations and frequencies where changes were induced for more than 70\% of participants.}}
\end{figure}

\subsection{Mold sets induced acoustic changes at distinct frequencies}\label{change_n}
To compare the acoustic effects between both mold sets, we mapped the number of participants for which each set induced marked changes in spectral amplitude across frequencies and elevations. Marked changes were defined as the mold-induced absolute difference in DTFs above a participant-specific threshold. The threshold was based on the variation between DTFs of participants' free ears and was defined as the mean spectral difference across elevations. The frequencies of marked changes within each set were similar across participants and differed between the two sets. For most participants, the first mold set induced changes to the central spectral notch in the 5–11 kHz band. The second set additionally affected spectral features in the neighboring 11-13 kHz band (Fig.~\ref{fig:FIG7} B-C). 

\begin{nofloatfigure}[b]
\centering
\includegraphics[width=8.5cm]{../figures/Figure8}
\caption{\label{fig:FIG8}{Vertical spectral information (VSI) and vertical spectral dissimilarity (VSI Dissimilarity) across pinna shapes. \textbf{A} Left and right ear spectral information in the 5.7–11.3 kHz band. Each marker represents one participant. Circles show spectral information measured with free ears, black and grey markers represent spectral information measured with earmolds 1 and 2, respectively. Relations between left and right ear spectral information are shown by regression lines. The relation between spectral information of participants' left and right ears persisted after mold insertion. \textbf{B} Spectral information with and without molds in the 5.7–11.3 kHz band. Both earmolds reduced spectral information in participants free ears. \textbf{C} Spectral dissimilarity measured across participants in the 5.7-13.5 kHz band. Grey circles show the spectral dissimilarity of free left and right ears between all combinations of two participants. Each colored marker shows the spectral dissimilarity induced by the molds for the left and right ears of one participant. The cloud of circles largely overlaps with the colored markers, indicating that acoustic differences induced by the molds were in the same range as inter-individual differences between participants’ free ears. \textbf{D} Spectral dissimilarity between the three pinna shapes in the 5.7-13.5 kHz band. The mean spectral dissimilarity between pinna shapes is indicated by the distance between markers ($\pm$ the standard error of the mean). Spectral dissimilarity between free ears and the first earmolds was similar to spectral dissimilarity between the first and second set of molds.}}
\end{nofloatfigure}

\subsection{Acoustic differences between consecutive pinna shapes were similar}
To test whether acoustic changes caused by the earmolds were in the same range, spectral dissimilarity between the different pinna shapes (free ears, mold 1 and mold 2) were computed in the 5.7–13.5 kHz frequency band. The dissimilarity between participants’ free ears and the first mold set was comparable to the dissimilarity between the first and the second set (see Fig.~\ref{fig:FIG8} D, spectral dissimilarity; free ears and mold 1: $W_{24} = 138, p = 0.363$). The second mold set induced larger changes to participants’ native ears than the first set (spectral dissimilarity; free ears and mold 1 vs. free ears and mold 2: $W_{24} = 138, p < .001$). To confirm that spectral changes induced by the molds were within a physiologically plausible range, spectral dissimilarities between free and modified ears of each participant were compared to dissimilarities between all possible pairs of participants’ free ears (Fig.~\ref{fig:FIG8} C). The overlap of distributions shows that spectral changes induced by the molds were comparable in magnitude to the natural spectrum of differences between individuals’ ears.

\begin{figure}[b]
\includegraphics[width=8.5cm]{../figures/Figure9}
\caption{\label{fig:FIG9}{Acoustic dissimilarity and localization performance. Relation between spectral dissimilarity and vertical error after mold insertion; \textbf{A} with mold 1 compared to free ears, \textbf{B} with mold 2 compared to free ears, \textbf{C} with mold 2 compared to mold 1. Behavior with the second earmolds was better explained by acoustic dissimilarity to the previously adapted molds than to participants’ native ears.}}
\end{figure}

\subsection{Spectral dissimilarity of earmolds did not fully explain acute effects on localization}
To test whether acoustic and behavioral effects of the earmolds were related, spectral dissimilarity between DTFs in the 5.7 – 13.5 kHz band with and without molds was compared to the reduction in participant’s localization performance after insertion of the earmolds. A weak non-significant correlation of increasing vertical error for larger acoustic differences was found for the first earmolds (Fig.~\ref{fig:FIG9} A, correlation of spectral dissimilarity between free ears and mold 1 and increase in vertical error after mold 1 insertion: $R = 0.43, p = 0.122$). No correlation was found for vertical localization performance and spectral dissimilarity between free ears and the second mold (Fig.~\ref{fig:FIG9} B). Because initial vertical localization accuracy with the second earmolds also depends on acoustic similarities to the previously learned set, differences in localization performance on the last day of adaptation to the first molds and the initial test with the second set were compared to the spectral dissimilarity between the both sets. A trend of increasing vertical error with larger acoustic dissimilarity between both mold sets can be seen in Fig.~\ref{fig:FIG9} C.

\section{\label{sec:4} Discussion}
Earmolds repeatedly disrupted vertical localization, but participants relearned sound localization with both mold sets within successive five-day adaptation periods. Individual adaptation performance did not improve with the second earmolds and learning a second set did not interfere with the previously learned representation. Localization accuracy with free ears slightly decreased throughout the experiment, whereas elevation perception remained unaffected. Spectral information of free ears correlated with behavior, and was similarly reduced by both mold sets. Acoustic dissimilarities between successive pinna shapes (free ears vs. the first molds and the first vs. the second molds) were similar, though each set induced changes at distinct frequencies.

\subsection{Localization with free ears} 
Two-dimensional localization performance with free ears was consistent with previous studies (\citet{hofman_relearning_1998}, \citet{makous_two-dimensional_1990}, \citet{oldfield_acuity_1984}, \citet{trapeau_adaptation_2015}). Vertical localization accuracy varied among individuals (\citet{wightman_headphone_1989}, \citet{populin_human_2008}), and correlated with individual spectral information in the 5.7-13.5 kHz range, suggesting that spectral features in this band contribute to vertical sound localization. However, individual performance may also depend on non-acoustic factors such as attention, perceptual abilities, and neural processes underlying spectral feature analysis (\citet{majdak_acoustic_2014}, \citet{andeol_sound_2013}). \citet{trapeau_fast_2016} showed that acoustic factors account for 25\% of individual differences in vertical localization accuracy. 

\subsection{Acoustic evaluation of the earmolds} 
Both sets of earmolds significantly reduced elevation perception and vertical accuracy, suggesting that the modified spectral-to-spatial mappings were sufficiently different from previously learned mappings to disrupt vertical localization. Horizontal accuracy and gain remained unaffected, indicating that binaural localization cues were not affected by the molds. To compare learning rates and adaptation persistence between the first and the second mold set, acoustic differences between successive pinna shapes should be similar. Prior research has shown that larger acoustic differences lead to stronger reduction in vertical localization performance and slower adaptation rates \citep{trapeau_fast_2016}. 
In this view, acoustic similarities between the first and the second set of earmolds could facilitate adaptation to the second set, because little cue reweighing would be required to accommodate to the new spectral mapping. Moreover, learning a second set of molds with very similar spectral cues could enhance the persistence of the previously learned mapping and counter possible effects of interference due to repeated learning of similar cues. 

Mold insertion altered spectral peaks and notches, which guide vertical localization (\citet{hofman_spectro-temporal_1998}, \citet{van_wanrooij_relearning_2005}, \citet{trapeau_fast_2016}). This was reflected by a reduction in spectral information at behaviorally relevant frequencies of participants' free ears. The volume and shape of the molds were adjusted to consistently disrupt elevation perception across participants. After adaptation to the first molds, the second set required larger silicone volumes to achieve a comparable reduction in elevation perception. The resulting smaller concha volumes interacted with shorter wavelengths, shifting the frequencies of spectral changes upwards. Earmolds might thus cause the formation of new spectral cues at higher frequencies \citep{van_wanrooij_relearning_2005}. This effect was observed for the first molds, which increased spectral information between 11.3 and 13.5 kHz compared to free ears. The effect of the second molds on localization was less pronounced for sounds at lower elevations. Apart from the outer ear, the head and torso also interact with incoming sound waves and provide coarse elevation cues for sources below the horizontal midline. Spectral cues created by the head and torso are situated at frequencies below 3 kHz (\citet{asano_role_1990}, \citet{algazi_elevation_2001}), and were likely unaffected by the molds. While spectral dissimilarity influenced localization performance, it did not fully account for individual differences, suggesting additional contributing factors \citep{trapeau_fast_2016}. Overall, both mold sets similarly reduced spectral information, with no significant differences in spectral dissimilarity between successive pinna shapes, enabling a direct comparison of behavioral effects across both mold sets.

\subsection{Adaptation to modified spectral cues} 
At the end of each adaptation period, vertical localization performance with earmolds trended toward performance with free ears.
However, elevation perception with molds remained poorer than the controls. Elevation gain plateaued at similar values as in previous studies (\citet{hofman_relearning_1998}, \citet{trapeau_fast_2016}, \citet{trapeau_encoding_2018}), suggesting a potential compression of auditory space along the vertical axis due to reduced spectral contrast (i.e., spectral information) provided by earmolds. However, as both mold sets similarly reduced spectral contrast, this did not confound adaptation comparisons. 

Learning rates and adaptation persistence exceeded those reported in (\citeyear{trapeau_fast_2016}), using an almost identical training paradigm. This study introduced gamification through a scoring system, allowing participants to compete with one another, and extended daily training sessions to 30 minutes of multiple short training runs interspersed with breaks (vs. a single 15-minute training run in the original paradigm). This may have increased motivation and attention, which are thought to be critical for perceptual learning (\citet{amitay_discrimination_2006}, \citet{mcgraw_introduction_2009}, \citet{molloy_less_2012}). Extensive training may introduce learning of artificial cues, such as loudspeaker characteristics. To address this, we equalized loudspeakers and confirmed the generalizability of cue relearning using stimuli of varying and unknown spectra, ensuring that adaptation was perceptual. However, the trade-off between experimental control and generalizeability to real world auditory scenes remains a limitation. 

Localization improvements varied among participants, ranging from minor adaptation to full recovery. Substantial inter-individual differences in spectral cue adaptation have been reported previously (\citet{hofman_relearning_1998}, \citet{van_wanrooij_relearning_2005}, \citet{trapeau_fast_2016}). However, the underlying causes of inter-individual differences in adaptation to spectral cues remain unclear. In \citeyear{trapeau_fast_2016}, acoustic factors explained 26\% of individual differences in localization performance after six days of adaptation and training. Non-acoustic factors, such as cognitive abilities and dispositional traits (e.g., sensitivity to reward) might explain individual variations in perceptual learning \citep{dale_individual_2021}. 

Earlier studies found that localization with free ears remained unaffected by adaptation to new spectral cues (\citet{hofman_relearning_1998}, \citet{van_wanrooij_relearning_2005}, \citet{carlile_relearning_2014}, \citet{trapeau_fast_2016}), indicating that the learned spectral-to-spatial mappings did not override pre-existing maps. In agreement with these results, adaptation persistence did not differ significantly between the successive earmolds, suggesting that adaptation to the second molds did not interfere with the previously learned representation. A slight increase in vertical error with the first molds after learning the second set may be due to fatigue after ten consecutive days of training and testing, causing participants to complete the localization tests more quickly. This could explain the reduced accuracy and precision with the second molds and the recovery of precision on day 15, after participants had a five-day break from the experiment. Similarly, the slight decrease in free-ear performance throughout the experiment may be attributed to fatigue rather than interference effects. 

Existing models of spectral location encoding posit that sound source elevation is estimated by comparing incoming sound spectra with stored spectral templates associated to different locations (\citet{hofman_spectro-temporal_1998}, \citet{langendijk_contribution_2002}). By the end of the second adaptation period, participants were able to access three different spectral mappings, which required the underlying neural populations to accommodate two additional spectral mappings in quick succession and with minimal interference. This raises the question at which stage of the auditory pathway the learned spectral templates are represented. Neural plasticity underlying spectral cue relearning may arise in subcortical structures. Neurons in the dorsal cochlear nucleus and their projections to the central nucleus of the inferior colliculus have been shown to encode spectral features in the cat \citep{davis_auditory_2003} and evidence exists for a topographical representation of auditory space in the superior colliculus of the ferret \citep{king_spatial_1987}. However, the ability to store multiple sets of spectral cues without interference suggests a surprisingly large capacity for this representation, which may be indicative of a cortical participation. 

Task repetition has been shown to increase learning rates in other cognitive domains. In the auditory pre-attentive domain, \citet{robinson_meta-adaptation_2016} demonstrated meta-adaptation in an intensity discrimination task, where cortical modulation of midbrain activity increased adaptation rates. However, our findings provide no evidence of meta-adaptation, as adaptation rate did not increase with repeated cue relearning. This may further support the notion of a cortical mechanism for spectral cue adaptation. In fact, adaptation rate was slightly lower in the second set. However, this might be attributed to acoustic factors, such as the increased spectral dissimilarity in the second set relative to participants’ native ears.

\subsection{Conclusion} 
Our results demonstrate that the auditory system can simultaneously encode three different maps of auditory space, indicating a surprisingly large capacity. The ability to store multiple spectral mappings without interference raises questions about the mechanisms underlying spectral cue encoding. While subcortical structures likely contribute to spectral processing, the large capacity is more suggestive of a cortical than subcortical mechanism. 

\bibliography{references.bib}

\newpage\textbf{Author Declarations}

The authors have no conflicts to disclose.
The present study was approved by the ethics advisory board of Leipzig University. Participants were informed about the relevant experimental procedures before providing their consent.

\\

\textbf{Data Availability}

The data that support the findings of this study are openly available at \href{https://github.com/pfriedrich-hub/HRTF_relearning_data/}{https://github.com/pfriedrich-hub/HRTF\_relearning\_data/}.


\end{document}
